% ===========================================================================
\section{Arithmetization summary}\label{sec:summary}
% ===========================================================================

We summarize the dimensions and constraint counts of the ECDSA
UAIR$^+$ above.

\begin{center}
\small
\begin{tabular}{ll}
\toprule
\textbf{Quantity} & \textbf{Value} \\
\midrule
Number of Shamir rounds            & $\mathrm{NUM\_SHAMIR\_ROUNDS} = 256$ \\
Final row index                    & $\mathrm{FINAL\_ROW} = 256$ \\
Active rows                        & $257$ \\
Trace length $n = 2^{\nu}$         & $2^{9} = 512$ \quad (smallest power of $2$ $\ge 257$) \\
Prime tuple $\mathbf{q}$           & $(p)$ \quad (single $\F_p$ branch; $p$ = $\secp$ base prime) \\
\midrule
Witness columns                    & $\mathbf{8}$ \quad (all native $\F_p$, no range checks) \\
Public-input columns               & $\mathbf{9}$ \\
Total trace columns                & $17$ \\
Lookup constraints                 & $0$ \\
Affine-combination lookups         & $0$ \\
\midrule
Constraint count                   & $\mathbf{11}$ \quad ($3$ doubling, $2$ add-scratch, $3$ output-selection, $3$ init-boundary) \\
Maximum constraint degree          & $\mathbf{6}$ \quad (after $\col{S\_ACTIVE}$ gating; \Cref{sec:degree}) \\
\bottomrule
\end{tabular}
\end{center}

\subsection{Why $8$ committed witness columns suffice}\label{sec:why-8}

A more direct arithmetization of the spec's Jacobian double-and-add
loop would commit, per row, separate columns for the doubled point
$P_{\mathrm{pa}}$, the addition output $P_{\mathrm{add}}$, the squared
intermediate $S = P_Y^2$, and the inversion-witness $Z_{\mathrm{inv}}$
for the final affine readout — totalling roughly $11$--$13$ committed
columns. The deployed UAIR shrinks this to $8$ via three deliberate
moves:
\begin{itemize}
\item \textbf{Inlined intermediates.}\
$S = P_Y^2$, $P_Y^4$, $P_{\mathrm{pa},Z}^2$, $P_{\mathrm{pa},Z}^3$,
$H^2$, $H^3$, and $P_{\mathrm{pa},X} H^2$ are all expanded inside the
constraint expressions rather than committed
(\Cref{rem:inlined}). Cost: max constraint degree rises from~$5$
to~$6$, but the global max is anyway~$6$ (from D4), so the bound is
unchanged.
\item \textbf{Conditional add via output-selection.}\
The addition output $P_{\mathrm{add}}$ is not committed; it is inlined
into the three output-selection constraints
\eqref{eq:next-x}--\eqref{eq:next-z}, which fold the conditional add
and the next-row chaining into one operation
(\Cref{rem:no-add-witness}). This removes three committed columns
that an explicit \emph{compute $P_{\mathrm{add}}$, then gate} approach
would need.
\item \textbf{Off-protocol affine readout.}\
The final-row Jacobian-to-affine conversion is performed by the
verifier off-protocol (\Cref{sec:boundary}). This removes the
$Z_{\mathrm{inv}}$ column and the affine output cells $(R_x, R_y)$,
saving up to three committed columns, and also removes the
non-infinity check $P_Z[256] \cdot Z_{\mathrm{inv}}[256] = 1$ and the
$x$-binding $P_X[256] \cdot Z_{\mathrm{inv}}[256]^2 = R_x$ from the
constraint set.
\end{itemize}
The retained witness columns
$(P_X, P_Y, P_Z, P_{\mathrm{pa},X}, P_{\mathrm{pa},Y}, P_{\mathrm{pa},Z}, H, R_a)$
are precisely those whose multiplicities in the constraint
expressions make inlining infeasible (each appears in at least two
constraint families).

\subsection{What the public input absorbs}\label{sec:public-absorbed}

Two further design moves push work out of the witness and into the
verifier's public-input construction. Both are noted earlier but
worth collecting:
\begin{itemize}
\item \textbf{$u_1, u_2$ and their bit-decompositions $b_1, b_2$
absorbed as public input.} The verifier computes
$u_1 = e \cdot s^{-1} \bmod n$ and $u_2 = r \cdot s^{-1} \bmod n$
off-protocol from the signature $(r, s)$ and message $e$, then
publishes $b_1, b_2$ along with off-protocol checks
$\sum_t b_i[t] \cdot 2^{256-t} = u_i$ in $\F_n$
(\Cref{rem:public-u}). This eliminates the entire $\F_n$-trace and
the bit-decomposition accumulator constraints; the loop reads the
bits directly via $\col{S\_ADD}$ and the per-row addend.

\item \textbf{Per-row addend coordinates absorbed as public input.}\
The verifier writes $T[t] = T_{b_1[t]\,b_2[t]} \in \{\OOO, G, Q, G+Q\}$
into the columns $(\col{PA\_X\_ADDEND}, \col{PA\_Y\_ADDEND})$ for each
row, with $\col{S\_ADD}[t] = \mathbf{1}\{T[t] \neq \OOO\}$
(\Cref{rem:addend}). This eliminates the four-way selector
constraints that an in-circuit \emph{compute the addend from the bits}
approach would need, and keeps the addition-formula constraints at
their stated degrees rather than pushing them up by the addend
selector.
\end{itemize}

\subsection{Composition with downstream UAIRs}

The retained $\col{S\_FINAL}$ selector (which is otherwise unused in
this UAIR) lets a downstream UAIR compose with the ECDSA UAIR by
adding final-row constraints — for example, an affine-conversion UAIR
that commits to $Z_{\mathrm{inv}}[256]$ and enforces
$P_Z[256] \cdot Z_{\mathrm{inv}}[256] = 1$ in-circuit, or a
gluing UAIR that binds the digest of an upstream SHA-256 UAIR to the
ECDSA message $e$ used to derive $u_1$. The implementation's
\texttt{EcdsaUair} + \texttt{Sha256CompressionSliceUair} side-by-side
merge follows exactly this pattern: the digest column
$\col{PA\_E}$ of the SHA half is read out-of-band by the verifier,
$u_1$ is computed off-protocol, and the resulting per-row
$(\col{PA\_X\_ADDEND}, \col{PA\_Y\_ADDEND})$ are written into the
ECDSA half's public input. The cross-binding is therefore
\emph{implicit}, enforced by the verifier supplying coherent public
inputs rather than by an in-circuit constraint, which is the same
composition pattern used by the SHA UAIR for its
$\col{PA\_A}/\col{PA\_E}$ digest exposure.
